\documentclass[12pt, a4paper]{article}

% Paquetes
\usepackage[utf8]{inputenc}
\usepackage[T1]{fontenc}
\usepackage[spanish]{babel}
\usepackage{amsmath, amssymb}
\usepackage{geometry}
\usepackage{booktabs}
\usepackage{array}
\usepackage{hyperref}
\usepackage{xcolor}
\usepackage{fancyhdr}
\usepackage{titlesec}
\usepackage{lmodern}
\usepackage{microtype}
\usepackage{tikz}
\usepackage{pgfplots}
\pgfplotsset{compat=1.18}
\usetikzlibrary{arrows.meta, decorations.pathmorphing, backgrounds, positioning}

\geometry{top=2.5cm, bottom=2.5cm, left=3cm, right=3cm}


\definecolor{azulms}{RGB}{25, 100, 180}
\definecolor{rojoquench}{RGB}{190, 40, 40}
\definecolor{verdedisco}{RGB}{30, 140, 80}
\definecolor{naranjabulge}{RGB}{210, 110, 20}
\definecolor{moradostarburst}{RGB}{120, 50, 160}
\definecolor{grisclaro}{RGB}{245, 245, 248}


\pagestyle{fancy}
\fancyhf{}
\rhead{\textcolor{azulms}{\small Contornos y validación de ruido}}
\lhead{\textcolor{azulms}{\small Pipeline ngVLA / VLA -- SINGS}}
\cfoot{\thepage}
\renewcommand{\headrulewidth}{0.5pt}


\titleformat{\section}{\large\bfseries\color{azulms}}{}{0em}{}[\titlerule]
\titleformat{\subsection}{\normalsize\bfseries\color{rojoquench}}{}{0em}{}



\begin{document}


% ── Portada
\begin{titlepage}
    \centering
    \vspace*{2cm}
    {\color{azulms}\rule{\linewidth}{1.5pt}}\\[0.4cm]
    {\LARGE\bfseries\color{azulms}
        Contornos de Detección\\[0.3cm]
        y Validación de Ruido}\\[0.4cm]
    {\color{azulms}\rule{\linewidth}{1.5pt}}\\[1.5cm]
    {\normalsize
        Pipeline de simulación ngVLA / VLA -- Galaxias SINGS}\\[2cm]
    {\normalsize Notas de dudas y respuestas}\\[0.5cm]
    {\normalsize \today}
    \vfill
\end{titlepage}

% ── Índice
\tableofcontents
\newpage

% ════════════════════════════════════════════════════════════════════════════
\section{Contexto: la petición del asesor}
% ════════════════════════════════════════════════════════════════════════════

Tras compartir la carpeta \texttt{collage/} (una imagen por galaxia con el mapa
de SFR y las cuatro configuraciones de telescopio -- ngVLA-A, ngVLA-B, VLA-A,
VLA-B -- en cada uno de los cinco redshifts simulados), el asesor pidió dos
cosas concretas:

\begin{enumerate}
    \item \textbf{Agregar contornos a las imágenes}, estimando el RMS de cada
          una y dibujando contornos a partir de $3\sigma$, para poder distinguir
          qué estructuras son emisión real de la galaxia y cuáles son
          simplemente ruido.
    \item \textbf{Medir el nivel de ruido} de cada imagen simulada y compararlo
          con el valor teórico de sensibilidad de la tabla de diseño del
          telescopio (ngVLA / VLA), como comprobación de que la simulación
          reproduce lo que el instrumento debería entregar.
\end{enumerate}

Este documento explica los conceptos estadísticos detrás de ambos puntos,
qué se encontró al implementarlos, y por qué el número de contornos visibles
cambia con el redshift.

% ════════════════════════════════════════════════════════════════════════════
\section{¿Qué es el RMS y qué significa \texorpdfstring{$\sigma$}{sigma}?}
% ════════════════════════════════════════════════════════════════════════════

En una imagen de radio sin fuente, el valor de cada píxel de fondo fluctúa
al azar alrededor de cero: es \textbf{ruido gaussiano}, producido por el
receptor, la electrónica y (en esta simulación) generado sintéticamente para
imitar esa estadística. El \textbf{RMS} (\textit{root mean square}) es la
dispersión característica de esas fluctuaciones:

\begin{equation}
    \mathrm{RMS} = \sigma = \sqrt{\left\langle (x - \bar{x})^2 \right\rangle},
    \label{eq:rms}
\end{equation}

es decir, "cuánto varían típicamente los píxeles de fondo alrededor de su
valor medio". Un valor de "$3\sigma$" es simplemente tres veces esa
dispersión.

\begin{center}
\begin{tikzpicture}
\begin{axis}[
    width=12cm, height=6.5cm,
    xlabel={Valor del píxel / $\sigma$},
    ylabel={Densidad de probabilidad},
    xmin=-5, xmax=5, ymin=0, ymax=0.45,
    grid=both, grid style={gray!20},
    axis lines=middle,
    title={\textbf{Distribución del ruido gaussiano y umbrales de detección}},
    title style={color=azulms},
    xtick={-5,-3,-2,-1,0,1,2,3,5},
    ytick=\empty,
    legend pos=north east,
    legend style={font=\footnotesize}
]

\addplot[thick, color=azulms, domain=-5:5, samples=200, fill=azulms, fill opacity=0.08]
    {1/sqrt(2*pi)*exp(-x^2/2)};
\addlegendentry{Ruido puro (gaussiano)}

\draw[dashed, color=rojoquench, thick] (axis cs:3,0) -- (axis cs:3,0.45);
\draw[dashed, color=rojoquench, thick] (axis cs:-3,0) -- (axis cs:-3,0.45);
\node[color=rojoquench, font=\small\bfseries] at (axis cs:3.6,0.40) {$3\sigma$};

\node[color=verdedisco, font=\small] at (axis cs:0,0.30) {$99.7\%$ del ruido};
\node[color=verdedisco, font=\small] at (axis cs:0,0.25) {cae dentro de $\pm 3\sigma$};

\end{axis}
\end{tikzpicture}
\end{center}

La razón de usar $3\sigma$ como umbral es puramente estadística: si el ruido
es gaussiano, solo un $0.3\%$ de los píxeles de puro ruido superarían ese
valor por azar. Por eso, si una región de la imagen supera $3\sigma$, es muy
poco probable que sea una fluctuación casual del ruido -- es una
\textbf{detección estadísticamente significativa}. Niveles más altos
($5\sigma$, $10\sigma$, $20\sigma$) marcan detecciones cada vez más seguras,
típicamente el núcleo de una fuente real.

En la práctica, el RMS no se calcula con la fórmula ingenua sobre toda la
imagen (eso incluiría la fuente y sesgaría el resultado hacia arriba), sino
con \textbf{sigma-clipping iterativo}
(\texttt{astropy.stats.sigma\_clipped\_stats}): se calcula la media y
desviación estándar, se descartan los píxeles que se alejan más de
$3\sigma$ de la media (probablemente señal, no fondo), y se repite el
cálculo solo con los píxeles restantes. Esto converge a una estimación
robusta del RMS del \textit{fondo}, sin que la propia fuente lo contamine.

% ════════════════════════════════════════════════════════════════════════════
\section{¿Qué significa el contorno en las imágenes ngVLA / VLA?}
% ════════════════════════════════════════════════════════════════════════════

El contorno (dibujado en color cian, \texttt{\#00ffcc}) traza el borde de las
regiones cuyo valor supera $3\sigma$, $5\sigma$, $10\sigma$ o $20\sigma$ del
RMS medido en esa misma imagen. Es clave entender lo que \textbf{no} hace: no
cambia los colores de la imagen, solo dibuja una línea sobre las zonas que
son estadísticamente significativas.

Esto es fundamental porque el color por sí solo \textbf{engaña al ojo}: con
mapas de color como \texttt{plasma} o \texttt{hot}, cualquier región que esté
cerca del percentil 99.5 de la imagen se ve "caliente" (amarilla o blanca),
sin importar si ese valor es una fuente real o simplemente el pico más alto
de una fluctuación de ruido correlacionado. El contorno es lo que permite
distinguir ambos casos:

\begin{itemize}
    \item \textbf{Con contorno alrededor}: la región supera el umbral de
          $3\sigma$ -- es una detección real.
    \item \textbf{Sin contorno, aunque se vea "caliente" por color}: es
          ruido correlacionado que por azar produce un pico visualmente
          similar, pero no supera el umbral estadístico.
\end{itemize}

En varias combinaciones de galaxia / redshift / configuración, el pico real
de la fuente está apenas por encima del ruido (o por debajo), así que el
contorno aparece como uno o ningún pequeño círculo aislado -- eso es un
resultado físico válido (no-detección o detección marginal), no un error de
la simulación.

% ════════════════════════════════════════════════════════════════════════════
\section{RMS medido vs. RMS teórico: el bug que se encontró}
% ════════════════════════════════════════════════════════════════════════════

\subsection{Qué es cada valor}

\begin{description}
    \item[\textcolor{azulms}{RMS teórico}] es la sensibilidad que promete
        la tabla de diseño del instrumento (p.\ ej.\ 28.44\,nJy/beam para
        ngVLA-A). Es lo que la simulación \textit{debería} reproducir si el
        ruido está bien generado.
    \item[\textcolor{rojoquench}{RMS medido}] es lo que efectivamente se
        calcula (con sigma-clipping) sobre cada imagen ya generada por el
        pipeline.
\end{description}

\subsection{El bug encontrado}

Al implementar esta comprobación, el RMS medido resultó \textbf{cinco órdenes
de magnitud menor} que el teórico ($\sim 7.8\times10^{-13}$ vs.\
$2.84\times10^{-8}$\,Jy/beam para ngVLA-A). La causa: la función
\texttt{add\_correlated\_noise} (duplicada en los ocho scripts de telescopio,
en \texttt{Modelo/} y \texttt{SINGS/}) atenuaba el ruido dos veces en
cascada:

\begin{enumerate}
    \item Primero dividía la desviación del ruido blanco por
          $\sqrt{N_{\mathrm{pix/beam}}}$, asumiendo que la convolución
          posterior "recuperaría" el nivel deseado.
    \item Pero el kernel gaussiano usado para correlacionar ese ruido se
          normalizaba \textbf{por su suma} (norma $L_1$, correcta para
          \textit{conservar flujo} al suavizar una imagen), en vez de por su
          norma $L_2$. Convolucionar ruido con un kernel $L_1$-normalizado
          \textbf{reduce su varianza} en un factor $\sqrt{\sum k_i^2} \ll 1$
          para un kernel ancho, en vez de preservarla.
\end{enumerate}

\subsection{La corrección}

\begin{itemize}
    \item El ruido blanco se genera directamente con
          $\mathrm{std} = \sigma_{\mathrm{beam}}$ (sin la predivisión).
    \item El kernel gaussiano se normaliza en norma $L_2$:
          \begin{equation}
              k_i \;\longrightarrow\; \frac{k_i}{\sqrt{\sum_j k_j^2}},
              \label{eq:kernel_l2}
          \end{equation}
          en vez de $k_i \to k_i / \sum_j k_j$.
\end{itemize}

Con esta normalización, convolucionar ruido de desviación
$\sigma_{\mathrm{beam}}$ con el kernel produce una imagen de ruido
correlacionado cuya desviación final es, de nuevo, $\sigma_{\mathrm{beam}}$
-- consistente con que cada píxel ya representa Jy/beam. Verificado
numéricamente antes de aplicar el cambio:

\begin{center}
\begin{tabular}{lcc}
\toprule
& \textbf{Antes (bug)} & \textbf{Después (corregido)} \\
\midrule
$\sigma$ tras convolución (ngVLA-A) & $6.9\times10^{-13}$\,Jy & $2.844\times10^{-8}$\,Jy \\
Valor teórico & \multicolumn{2}{c}{$2.844\times10^{-8}$\,Jy} \\
\bottomrule
\end{tabular}
\end{center}

Tras corregir y regenerar las imágenes, el RMS medido quedó entre
$0.7\times$ y $1.0\times$ el valor teórico -- el rango esperado, ya que cada
imagen es una única realización aleatoria del ruido, no un promedio
infinito de muchas.

% ════════════════════════════════════════════════════════════════════════════
\section{¿Por qué aparecen más contornos a redshift alto?}
% ════════════════════════════════════════════════════════════════════════════

Al comparar los collages se observó que a $z=0.4$ casi no hay contornos
(cero o uno), mientras que a $z=3.0$--$5.0$ aparecen varios. Se verificó
contando "blobs" (regiones conectadas sobre $3\sigma$) promediando las 14
galaxias disponibles:

\begin{center}
\begin{tabular}{cccc}
\toprule
$z$ & \textbf{SFR asignado} & \textbf{Blobs promedio} & \textbf{Área $>3\sigma$} \\
    & [$M_\odot$/yr] & (ngVLA-A) & \\
\midrule
0.4 & 1.5  & 1.8 & 0.100\% \\
1.0 & 7    & 1.1 & 0.110\% \\
2.0 & 40   & 0.8 & 0.105\% \\
3.0 & 100  & 1.2 & 0.189\% \\
5.0 & 300  & 2.3 & 0.182\% \\
\bottomrule
\end{tabular}
\end{center}

La tendencia es real y se explica por dos efectos combinados:

\subsection{1. La tabla de SFR hace la galaxia intrínsecamente más luminosa a alto \texorpdfstring{$z$}{z}}

La tabla de redshift--SFR (\texttt{SINGS/utils.py}, basada en
Leslie et al.\ 2020 para $M_\star = 10^{10}\,M_\odot$) no es arbitraria:
representa el SFR \textit{típico} de una galaxia de masa fija en la
Secuencia Principal de formación estelar en cada época. Esa secuencia sube
fuertemente con $z$ -- el universo joven formaba estrellas mucho más rápido.
Como la luminosidad radio es directamente proporcional al SFR
($L = \mathrm{SFR}/4.87\times10^{-29}$), a $z=5$ la fuente es
\textbf{$\sim$200 veces} más luminosa intrínsecamente que a $z=0.4$.

\subsection{2. Esa ganancia compite con el oscurecimiento por distancia, y casi lo compensa}

Alejar la fuente reduce el flujo observado como $\sim 1/D_L^2$, y la
distancia de luminosidad $D_L$ crece muchísimo con $z$. Sin embargo, al medir
el flujo total (sin ruido) de una misma galaxia (NGC\,628) en los cinco
redshifts, el resultado es sorprendentemente parecido en todos los casos:

\begin{center}
\begin{tabular}{cc}
\toprule
$z$ & \textbf{Flujo total} [Jy] \\
\midrule
0.4 & $1.52\times10^{-6}$ \\
1.0 & $8.56\times10^{-7}$ \\
2.0 & $9.98\times10^{-7}$ \\
3.0 & $1.02\times10^{-6}$ \\
5.0 & $1.02\times10^{-6}$ \\
\bottomrule
\end{tabular}
\end{center}

Es decir: la tabla está construida de forma que el aumento de SFR compensa
\textit{aproximadamente} el oscurecimiento cosmológico -- pero no
exactamente. Además, el \textbf{ruido del instrumento no depende de $z$}
(28.44\,nJy/beam para ngVLA-A, siempre el mismo, sin importar qué tan lejos
esté la fuente). Entonces, si el flujo se mantiene similar o sube levemente
con $z$ mientras el ruido queda fijo, la relación señal/ruido mejora
ligeramente hacia $z$ alto, y más píxeles cruzan el umbral fijo de
$3\sigma$.

\subsection{En resumen}

A $z=0.4$ la galaxia es intrínsecamente muy débil en radio (SFR $=1.5\,
M_\odot$/yr es poco), así que aunque está cerca, casi no hay señal por
encima del ruido. A $z$ alto, aunque está mucho más lejos, se le asignó un
SFR mucho mayor -- como corresponde a una galaxia típica de esa época -- y
esa ganancia de luminosidad gana la carrera contra la distancia. Por eso se
ven más contornos (más detecciones) hacia redshift alto.

% ════════════════════════════════════════════════════════════════════════════
\section{Resumen de cambios aplicados al pipeline}
% ════════════════════════════════════════════════════════════════════════════

\begin{enumerate}
    \item \textbf{Corrección del bug de ruido} en \texttt{add\_correlated\_noise}
          (los cuatro scripts de telescopio en \texttt{SINGS/}): generación
          directa del ruido blanco con $\mathrm{std}=\sigma_{\mathrm{beam}}$
          y normalización $L_2$ del kernel gaussiano.
    \item \textbf{\texttt{collage.py}}: se agregó \texttt{estimate\_rms()}
          (sigma-clipping) y la anotación de RMS medido vs.\ teórico bajo
          cada panel de telescopio.
    \item Los cuatro paneles de telescopio se redibujan con
          \texttt{aplpy.FITSFigure}, incluyendo contornos WCS-aware
          (\texttt{show\_contour}) desde $3\sigma$, y el beam real leído del
          header FITS (\texttt{add\_beam}).
    \item Color del contorno cambiado de blanco a cian
          (\texttt{\#00ffcc}) para que contraste con los mapas de color
          \texttt{plasma} y \texttt{hot}, que terminan en blanco/amarillo
          pálido justo donde caen los picos de $3\sigma$.
    \item El panel de "Mapa radio" (SFR sin instrumento, antes de la
          convolución con el beam) se dejó \textbf{sin contorno}: ese mapa
          tiene ruido pixel a pixel sin suavizar y un posible gradiente de
          fondo heredado del H$\alpha$ de entrada, por lo que un contorno de
          silueta limpio requeriría trabajo adicional de suavizado y
          caracterización del fondo -- un problema de calidad de datos
          aparte del propósito original de esta tarea.
\end{enumerate}

\end{document}
